O Sr. Jones, proprietário da Granja do Solar, fechou o galinheiro à noite, mas estava bêbado demais para lembrar-se de fechar também as vigias. Com o facho de luz da sua lanterna balançando de um lado para o outro, atravessou cambaleante o pátio, tirou as botas na porta dos fundos, tomou um último copo de cerveja do barril que havia na copa, e foi para a cama, onde sua mulher já ressonava.

Tão logo apagou-se a luz do quarto, houve um grande alvoroço em todos os galpões da granja. Correra, durante o dia, o boato de que o velho Major, um porco que já se sagrara grande campeão numa exposição, tivera um sonho muito estranho na noite anterior e desejava contá-lo aos outros animais. Haviam combinado encontrar-se no celeiro, assim que Jones se retirasse. O velho Major (chamavam-no assim, muito embora ele houvesse comparecido à exposição com o nome de ``Beleza de Willingdon'') gozava de tão alto conceito na granja, que todos estavam dispostos a perder uma hora de sono só para ouvi-lo.

Ao fundo do grande celeiro, sobre uma espécie de estrado, estava o Major refestelado em sua cama de palha, sob um lampião que pendia de uma viga. Com doze anos de idade, já bastante corpulento, era ainda um porco de porte majestoso, com um ar sábio e benevolente, a despeito de suas presas jamais terem sido cortadas. Os outros animais chegavam e punham-se à vontade, cada qual a seu modo. Os primeiros foram os três cachorros, Ferrabrás, Lulu e Cata-vento, depois os porcos, que se sentaram sobre a palha, em frente ao estrado. As galinhas empoleiraram-se nas janelas, as pombas voaram para os caibros do telhado, as ovelhas e as vacas deitaram-se atrás dos porcos e ali ficaram a ruminar. Os dois cavalos de tração, Sansão e Quitéria, chegaram juntos, andando lentamente e pousando no chão os enormes cascos peludos, com grande cuidado para não machucar qualquer animalzinho porventura oculto na palha. Quitéria era uma égua volumosa, matronal, já chegada à meia-idade, cuja silhueta não mais se recompusera após o nascimento do quarto potrinho. Sansão era um bicho enorme, de quase um metro e noventa de altura, forte como dois cavalos. A mancha branca do focinho dava-lhe um certo ar de estupidez e, realmente, não tinha lá uma inteligência de primeira ordem, embora fosse grandemente respeitado pela retidão de caráter e pela tremenda capacidade de trabalho. Depois dos cavalos chegaram Maricota, a cabra branca, e Benjamim, o burro. Benjamim era o animal mais idoso da fazenda, e o mais moderado. Raras vezes falava e, normalmente, quando o fazia, era para emitir uma observação cínica --- para dizer, por exemplo, que Deus lhe dera uma cauda para espantar as moscas e que, no entanto, seria mais do seu agrado não ter nem a cauda nem as moscas. Era o único dos animais que nunca ria. Quando lhe perguntavam por que, respondia não ver motivo para riso. Não obstante, sem que o admitisse abertamente, tinha certa afeição por Sansão; normalmente passavam os domingos juntos no pequeno potreiro existente atrás do pomar, pastando lado a lado em silêncio.

\begin{figure}[h]
	\vspace{-15mm}
	\centering % Centralizar
	\includegraphics[scale=0.3]{img/CAP 01} % arquivo da imagem.
	\caption{DISCURSO DE MAJOR}
    end{figure} 
  
Mal se haviam acomodado os dois cavalos quando uma ninhada de patinhos órfãos desfilou celeiro adentro, piando baixinho e procurando um lugar onde não fossem pisoteados. Quitéria protegeu-os com a pata dianteira e os patinhos ali se aconchegaram, caindo no sono. No último instante, Mimosa, a égua branca, vaidosa e fútil, que puxava a aranha do Sr. Jones, entrou, requebrando-se graciosamente e chupando um torrão de açúcar. Tomou um lugar bem à frente e ficou meneando a sua crina branca, na esperança de chamar atenção para as fitas vermelhas que a adornavam. Finalmente, chegou o gato, que procurou, como sempre, o lugar mais morno, enfiando-se entre Sansão e Quitéria; ressonou satisfeito durante toda a fala do Major, sem ouvir uma só palavra.

Todos os animais estavam presentes, exceto Moisés, o corvo domesticado, que dormia fora, num poleiro junto à porta dos fundos. Quando o Major os viu bem acomodados e aguardando atentamente, limpou a garganta e começou:

--- ``Camaradas, já ouvistes, por certo, algo a respeito do estranho sonho que tive na noite passada. Entretanto, falarei do sonho mais tarde. Antes, há coisas a dizer. Sei, camaradas, que não estarei convosco por muito tempo e antes de morrer considero uma obrigação transmitir-vos o que tenho aprendido sobre o mundo. Já vivi bastante e muito tenho refletido na solidão da minha pocilga. Creio poder afirmar que compreendo a natureza da vida sobre esta terra, tão bem quanto qualquer outro animal. É sobre isso que desejo falar-vos.

``Então, camaradas, qual é a natureza da nossa vida? Enfrentemos a realidade: nossa vida é miserável, trabalhosa e curta. Nascemos, recebemos o mínimo de alimento necessário para continuar respirando e os que podem trabalhar são forçados a fazê-lo até a última parcela de suas forças; no instante em que nossa utilidade acaba, trucidam-nos com hedionda crueldade. Nenhum animal, na Inglaterra, sabe o que é felicidade ou lazer, após completar um ano de vida. Nenhum animal, na Inglaterra, é livre. A vida de um animal é feita de miséria e escravidão: essa é a verdade nua e crua.

``Será isso, apenas, a ordem natural das coisas? Será esta nossa terra tão pobre que não ofereça condições de vida decente aos seus habitantes? Não, camaradas, mil vezes não! O solo da Inglaterra é fértil, o clima é bom, ela pode oferecer alimentos em abundância a um número de animais muitíssimo maior do que o existente. Só esta nossa fazenda comportaria uma dúzia de cavalos, umas vinte vacas, centenas de ovelhas --- vivendo todos com uma dignidade que, agora, está além de nossa imaginação. Por que, então, permanecemos nesta miséria? Porque quase todo o produto do nosso esforço nos é roubado pelos seres humanos. Eis aí, camaradas, a resposta a todos os nossos problemas. Resume-se em uma só palavra: Homem. O homem é o nosso verdadeiro e único inimigo. Retire-se da cena o Homem, e a causa principal da fome e da sobrecarga de trabalho desaparecerá para sempre.

``O Homem é a única criatura que consome sem produzir. Não dá leite, não põe ovos, é fraco demais para puxar o arado, não corre o suficiente para alcançar uma lebre. Mesmo assim, é o senhor de todos os animais. Põe-nos a trabalhar, dá-nos de volta o mínimo para evitar a inanição e fica com o restante. Nosso trabalho amanha o solo, nosso estrume o fertiliza e, no entanto, nenhum de nós possui mais do que a própria pele. As vacas, que aqui vejo à minha frente, quantos litros de leite terão produzido este ano? E que aconteceu a esse leite, que deveria estar alimentando robustos bezerrinhos? Desceu pela garganta dos nossos inimigos. E as galinhas, quantos ovos puseram este ano, e quantos se transformaram em pintinhos? Os restantes foram para o mercado, fazer dinheiro para Jones e seus homens. E você, Quitéria, diga-me onde estão os quatro potrinhos que deveriam ser o apoio e o prazer da sua velhice? Foram vendidos com a idade de um ano --- nunca você tornará a vê-los. Como paga pelos seus quatro partos e por todo o seu trabalho no campo, que recebeu você, além de ração e baia?

``Mesmo miserável como é, nossa vida não chega ao fim de modo natural. Não me queixo por mim, que tive até muita sorte. Estou com doze anos e sou pai de mais de quatrocentos porcos. Isto é a vida normal de um varrão. Mas, no fim, nenhum animal escapa ao cutelo. Vós, jovens leitões que estais sentados à minha frente, não escapareis de guinchar no cepo dentro de um ano. Todos chegaremos a esse horror, as vacas, os porcos, as galinhas, as ovelhas, todos. Nem mesmo os cavalos e os cachorros escapam a esse destino. Você, Sansão, no dia em que seus músculos fortes perderem a rigidez, Jones o mandará para o carniceiro e você será degolado e fervido para os cães de caça. Quanto aos cachorros, depois de velhos e desdentados, Jones amarra-lhes uma pedra ao pescoço e joga-os na primeira lagoa.

``E lembrai-vos, camaradas, jamais deixai fraquejar vossa decisão. Nenhum argumento poderá deter-vos. Fechai os ouvidos quando vos disserem que o Homem e os animais têm interesses comuns, que a prosperidade de um é a prosperidade dos outros. É tudo mentira. O Homem não busca interesses que não os dele próprio. Que haja entre nós uma perfeita unidade, uma perfeita camaradagem na luta. Todos os homens são inimigos, todos os animais são camaradas.''

Nesse momento houve uma tremenda confusão. Enquanto o Major falava, quatro ratos haviam emergido de seus buracos e estavam sentados nas patinhas de trás, a ouvi-lo. De repente, os cachorros lhes deram pela presença, e somente devido à rapidez com que sumiram nos buracos foi que os ratos conseguiram escapar com vida. O Major levantou a pata, pedindo silêncio.

--- ``Camaradas --- disse ele ---, eis aí um ponto que precisa ser esclarecido. As criaturas selvagens, tais como os ratos e os coelhos, serão nossos amigos ou nossos inimigos? Coloquemos o assunto em votação. Apresento à assembleia a seguinte questão: os ratos são camaradas?''

A votação foi realizada imediatamente e concluiu-se, por esmagadora maioria, que os ratos eram camaradas. Houve apenas quatro votos contra, dos três cachorros e do gato que, depois se descobriu, votara pelos dois lados. O Major prosseguiu:

--- ``Pouco mais tenho a dizer. Repito apenas: lembrai-vos sempre do vosso dever de inimizade para com o Homem e todos os seus desígnios. Qualquer coisa que ande sobre duas pernas é inimigo, qualquer coisa que ande sobre quatro pernas, ou tenha asas, é amigo. Lembrai-vos também de que na luta contra o Homem não devemos assemelhar-nos a ele. Mesmo quando o tenhais derrotado, evitai seus vícios. Animal nenhum deve morar em casas, nem dormir em camas, nem usar roupas, nem beber álcool, nem fumar, nem tocar em dinheiro, nem fazer comércio. Todos os hábitos do Homem são maus. E, principalmente, jamais um animal deverá tiranizar outros animais. Todos os animais são iguais.

``E agora, camaradas, vou contar-vos o sonho que tive na noite passada. Não sei como explicá-lo. Foi um sonho sobre como será o mundo quando o Homem desaparecer. Mas lembrou-me algo que há muito eu esquecera. Há anos, quando eu ainda era um leitãozinho, minha mãe e as outras porcas costumavam cantar uma antiga canção da qual só conheciam a melodia e as três primeiras palavras. Na minha infância aprendi a melodia, depois a esqueci. Na noite passada, entretanto, ela me voltou à memória. O mais interessante é que me lembrei também dos versos --- os quais, tenho certeza, foram cantados pelos animais de antanho, e depois esquecidos durante várias gerações. Vou cantar essa canção, camaradas. Estou velho e minha voz é rouca, mas quando vos houver ensinado a melodia, podereis cantá-la melhor do que eu. Chama-se \textit{Bichos da Inglaterra}.''

O velho Major limpou a garganta e começou a cantar. De fato, a voz era roufenha, mas ele cantava razoavelmente, e a melodia era bem movimentada, algo entre \textit{Clementine} e \textit{La Cucaracha}. Os versos diziam o seguinte:

\begin{verse}
Bichos ingleses e irlandeses, \\
Bichos de todas as partes! \\
Eis a mensagem de esperança, \\
No futuro que virá!

Cedo ou tarde virá o dia, \\
Cairá a tirania \\
E os campos todos da Inglaterra \\
Só aos bichos caberão!

Não mais argolas em nossas ventas, \\
Dorsos livres dos arreios, \\
Freios e esporas, descartados, \\
Chicotadas abolidas!

Muito mais ricos do que sonhamos \\
Possuiremos daí por diante \\
O trigo, o feno, e a cevada, \\
Pasto, aveia e feijão!

Brilham os campos da Inglaterra, \\
Águas puras rolarão. \\
Ventos leves soprarão \\
Saudando a redenção!

Lutemos todos por esse dia \\
Mesmo que nos custe a vida! \\
Cavalos, vacas, perus e gansos, \\
Liberdade conquistemos!

Bichos ingleses e irlandeses, \\
Bichos de todas as partes! \\
No futuro que virá!
\end{verse}

O canto levou os animais à mais extrema excitação. Antes de o Major chegar ao fim, já haviam começado a cantar por conta própria. Até os mais estúpidos pegaram a melodia e algumas palavras; os mais espertos, como os porcos e os cachorros, decoraram a canção em poucos minutos. Então, depois de alguns ensaios preliminares, toda a granja atacou \textit{Bichos da Inglaterra}, em formidável uníssono. As vacas mugiam a canção, os cachorros latiam-na, as ovelhas baliam-na, os cavalos relinchavam-na, os patos grasnavam-na. Tal foi o enlevo, que cantaram de ponta a ponta, cinco vezes sucessivamente, e teriam continuado a noite inteira se não fossem interrompidos.

Infelizmente, o alarido acordou Jones, que pulou da cama certo de que havia raposa no pátio. Deu de mão na espingarda, sempre pronta a um canto do quarto, e descarregou-a na escuridão. O chumbo foi encravar-se na parede do celeiro, e a reunião dispersou-se num abrir e fechar de olhos. Cada qual correu para seu pouso. As aves saltaram para os poleiros, o gado deitou-se na palha e, em poucos instantes, toda a fazenda dormia.